\clearpage
\onecolumn
\section{Qualitative Gallery}
\label{sec:gallery}

The following pages collect the clearest qualitative results we have rendered.
They are presentation views, not additional samples or a hidden selection stage:
every crop appears in the committed locked-16 or fixed PHerc1447 review corpus.

Both galleries show the earlier v31 student at $t=0.45$ rather than the
checkpoint reported in Table~\ref{tab:headline}, and have not been regenerated
against the shipped weights; see Sec.~\ref{sec:figure-provenance}.
They are included because the failure modes they illustrate---longitudinal
growth on locked slices, and de-blob behaviour on blind cubes---are the ones the
method targets, and because the released M7 baseline shown beside every panel is
unchanged.

\begin{center}
  \figureasset{gallery_locked_growth}{0.86\textwidth}{7.4in}
  \captionsetup{hypcap=false}
  \captionof{figure}{Locked PHerc0139 growth gallery. Every registered crop is
  shown as CT, released M7, projected fine-teacher target, and the raw student.
  These examples emphasize learned longitudinal regularity while keeping the
  released state of the art visible.}
  \label{fig:gallery-locked}
\end{center}

\clearpage
\begin{center}
  \figureasset{gallery_blind_antiblob}{0.86\textwidth}{7.4in}
  \captionsetup{hypcap=false}
  \captionof{figure}{Blind PHerc1447 anti-blob gallery. The released M7
  foreground is cyan and the raw student's foreground is magenta, each on the
  identical CT slice. No fine teacher exists for these views.}
  \label{fig:gallery-blind}
\end{center}

\clearpage
\twocolumn
